% ============================================================================
\clearpage
\section{De la información implícita Q a pronósticos reales P}\label{sec:pronostico}
% ============================================================================

\ficha{Variables de $\Q$ (factores, colas, innovaciones), información base y rendimientos realizados
después de cada snapshot.}
{Entrenar modelos supervisados por objetivo y horizonte, corregir cruces de cuantiles y combinar.}
{Distribuciones físicas $\P$ por horizonte: cuantiles, probabilidades de evento y volatilidad
esperada.}
{Etiquetas alineadas con el instante de decisión; mejora incremental medida dentro de una misma
familia de modelos.}

\subsection{Por qué Q no es P}\label{sec:qp}

$\Q$ y $\P$ son dos distribuciones distintas del mismo rendimiento. $\Q$ es la que hace coherentes los
precios de las opciones; $\P$ describe con qué frecuencia ocurren las cosas. Las une un núcleo de
valoración (factor estocástico de descuento proyectado) $m$:
\begin{equation}
q(x)=\frac{m(x)\,p(x)}{\E^{\P}[m]}.
\label{eq:nucleo}
\end{equation}
Sin supuestos adicionales $m$ no está identificado: la misma $q$ es compatible con muchas parejas $(m,p)$.
Los intentos de ``recuperar'' $\P$ a partir de $\Q$ requieren supuestos fuertes \citep{ross2015} cuya
validez está en discusión \citep{borovicka2016}. Para índices bursátiles se documenta, en promedio, una
prima de varianza: la dispersión implícita tiende a superar a la realizada \citep{carrwu2009}. La
\figref{fig:q_vs_p} ilustra un supuesto de ese tipo y la \figref{fig:nucleo} el núcleo que implicaría.

\grafica[0.62]{s4_q_vs_p}{Una distribución $\Q$ y una $\P$ supuesta para el mismo activo. La media de $\Q$
está fijada por el forward; la de $\P$ incorpora una prima. La cola izquierda de $\Q$ es más pesada.
Ilustrativo: la relación real entre ambas es lo que se estima, no lo que se supone.}{fig:q_vs_p}

\grafica[0.62]{s4_nucleo}{Cociente $q/p$ implícito en el supuesto de la \figref{fig:q_vs_p}. Otra elección
de $p$ daría otro cociente con la misma $q$: por eso $\P$ se estima con rendimientos realizados y no se
deduce de $\Q$.}{fig:nucleo}

\subsection{Diseño supervisado y alineación temporal}\label{sec:alineacion}

Los modelos se entrenan con rendimientos efectivamente realizados \emph{después} de cada snapshot. Las
variables de $\Q$ entran como regresores; no se etiquetan como probabilidades físicas por decreto. La
etiqueta de horizonte $h$ solo existe en $t+h$, y hasta entonces no puede usarse para entrenar, calibrar ni
actualizar intervalos (\figref{fig:alineacion}).

\begin{figure}[H]
\centering
\resizebox{\linewidth}{!}{% Alineación temporal del aprendizaje supervisado: variables, decisión y etiqueta.
\begin{tikzpicture}[x=0.62cm,y=1cm]
% Ventana de variables (pasado)
\fill[qazul!14] (-4,0.3) rectangle (0,1.9);
\node[nota,text=tinta,align=center] at (-2,1.1) {\textbf{variables}\\ información hasta\\ el corte};
% Instante de decisión
\draw[tinta,line width=0.9pt] (0,2.35) -- (0,-0.1);
\node[nota,anchor=south,text=tinta] at (0,2.35) {\textbf{decisión en $t$}};
% Etiquetas por horizonte
\foreach \y/\h in {1.55/1,1.0/5,0.45/20}{
  \fill[qnaranja!25] (0,\y-0.15) rectangle (\h,\y+0.15);
  \draw[qnaranja,line width=0.7pt] (\h,\y-0.15) -- (\h,\y+0.15);
  \node[nota,anchor=west,text=tinta] at (\h+0.2,\y) {$h=\h$};}
% Eje de sesiones
\draw[guia,line width=0.7pt] (-4,0) -- (22,0);
\foreach \s in {0,1,...,21} \draw[guia] (\s,0) -- (\s,-0.08);
\foreach \s/\e in {0/t,1/t+1,5/t+5,20/t+20} \node[nota,anchor=north] at (\s,-0.1) {$\e$};
\node[nota,anchor=north,text width=12cm,align=center] at (9,-0.5) {La etiqueta de horizonte $h$ es el rendimiento
  en $(t,\,t+h]$. Solo después de madurar, en $t+h$, puede entrar al entrenamiento, a la calibración o a la
  actualización conformal.};
\end{tikzpicture}
}
\caption{Alineación temporal. Las variables se calculan con información hasta el corte; la etiqueta es el
rendimiento posterior al instante de decisión y madura al final de su horizonte.}
\label{fig:alineacion}
\end{figure}

\subsection{Escalera de modelos}\label{sec:escalera}

El orden propuesto va de lo simple a lo flexible, y cada peldaño debe ganarle al anterior fuera de muestra
para justificar su complejidad (\figref{fig:escalera}):

\begin{enumerate}
\item \textbf{Referencia histórica con volatilidad dinámica}: cuantiles de residuos estandarizados
históricos escalados por una volatilidad EWMA o GARCH.
\item \textbf{Regresión cuantílica regularizada} \citep{koenker1978}: para cada nivel $\tau$,
\begin{equation}
\hat\beta_\tau=\arg\min_\beta\sum_t\rho_\tau\bigl(y_t-z_t^\top\beta\bigr)+\lambda\lVert\beta_{-0}\rVert_1,
\qquad \rho_\tau(e)=e\,\bigl(\tau-\mathbb 1\{e<0\}\bigr),
\label{eq:qr}
\end{equation}
un programa lineal (\figref{fig:pinball}).
\item \textbf{Boosting cuantílico} con la misma pérdida y corrección de cruces.
\item \textbf{NGBoost} \citep{duan2020}, que estima con gradiente natural los parámetros de una familia de
distribuciones condicional. La familia y sus colas se eligen por validación, sin imponer normalidad por
comodidad.
\end{enumerate}

\begin{figure}[H]
\centering
\resizebox{\linewidth}{!}{% Escalera de modelos P y combinación de pronósticos.
\begin{tikzpicture}[
  m/.style={caja=qnaranja,text width=3.05cm,minimum height=16mm},
  node distance=5mm and 5mm]
\node[m] (m1) {\textbf{Referencia histórica}\\ volatilidad dinámica y residuos estandarizados};
\node[m,right=of m1] (m2) {\textbf{Regresión cuantílica}\\ regularizada ($L_1$), un modelo por $\tau$};
\node[m,right=of m2] (m3) {\textbf{Boosting cuantílico}\\ con corrección de cruces};
\node[m,right=of m3] (m4) {\textbf{NGBoost}\\ distribución condicional; familia elegida en validación};
\draw[flecha] (m1) -- (m2); \draw[flecha] (m2) -- (m3); \draw[flecha] (m3) -- (m4);
\node[nota,above=1mm of m2.north east,anchor=south] {cada peldaño debe ganarle al anterior fuera de muestra};
\coordinate (medio) at ($(m2.south)!0.5!(m3.south)$);
\node[caja=qnaranja,text width=7.2cm,minimum height=12mm,below=11mm of medio] (comb)
  {\textbf{Combinación}: pesos $w_k\ge 0$, $\sum_k w_k=1$, aprendidos en validación temporal
  y comparados con pesos iguales};
\coordinate (bus) at ($(comb.north)+(0,0.45)$);
\foreach \m in {m1,m2,m3,m4} \draw[line width=0.6pt,draw=tinta2] (\m.south) |- (bus);
\draw[flecha] (bus) -- (comb.north);
\node[caja=tenue,text width=13.9cm,minimum height=9mm,below=6mm of comb] (obj)
  {\textbf{Un objetivo explícito por horizonte} (1, 5 y 20 sesiones): rendimiento $\cdot$ probabilidad de superar
  costos $\cdot$ cuantiles de pérdida $\cdot$ volatilidad futura};
\draw[flecha] (comb) -- (obj);
\end{tikzpicture}
}
\caption{Escalera de modelos $\P$, combinación y objetivos por horizonte.}
\label{fig:escalera}
\end{figure}

\grafica[0.56]{s4_pinball}{Pérdida cuantílica para $\tau=0.05$, $0.50$ y $0.95$. Su asimetría es la que
hace que el minimizador sea el cuantil $\tau$.}{fig:pinball}

Estimar cada cuantil por separado puede producir \emph{cruces}: en una zona de extrapolación el cuantil
del 25\% queda por debajo del 5\% (\figref{fig:cruces}). El reordenamiento creciente de los cuantiles
estimados en cada punto corrige el problema sin empeorar la estimación
\citep{chernozhukov2010} (\figref{fig:reordenado}).

\grafica[0.6]{s4_cruces}{Regresiones cuantílicas estimadas por separado con 70 observaciones sintéticas. Fuera del rango observado
de la variable, dos cuantiles se cruzan (punto rojo).}{fig:cruces}

\grafica[0.6]{s4_reordenado}{Los mismos cuantiles después de reordenarlos en cada valor de la variable: siempre crecientes en
$\tau$.}{fig:reordenado}

\subsection{Objetivos por horizonte}\label{sec:objetivos}

Cada objetivo tiene su pérdida y se entrena por separado. La información de opciones podría mejorar el
riesgo sin mejorar el signo del rendimiento; ambas conclusiones son útiles y se muestran por separado.

\begin{table}[H]
\centering\small
\caption{Objetivos, pronósticos y pérdidas. Horizontes candidatos: 1, 5 y 20 sesiones.}
\label{tab:objetivos}
\begin{tabularx}{\linewidth}{@{}l Y Y@{}}
\toprule
\encab{Objetivo} & \encab{Pronóstico} & \encab{Pérdida y diagnóstico}\\
\midrule
Rendimiento & Distribución completa; mediana & CRPS; pinball en $\tau=0.5$; calibración de cuantiles\\
Superar costos & Probabilidad de que el rendimiento neto sea positivo & Brier y log score; diagrama de fiabilidad\\
Cuantiles de pérdida & $q_{0.05}$ y $q_{0.10}$ & Pinball; cobertura y ancho de intervalos\\
Volatilidad futura & Varianza realizada esperada & QLIKE, robusta a proxies ruidosos \citep{patton2011}\\
\bottomrule
\end{tabularx}
\end{table}

\subsection{Combinación de pronósticos}\label{sec:combinacion}

Los pronósticos se combinan con pesos no negativos que suman uno, aprendidos en validación temporal
---por ejemplo, minimizando el CRPS de la mezcla o la pérdida cuantílica del promedio de cuantiles--- y se
comparan siempre con pesos iguales, que en la práctica son difíciles de superar \citep{timmermann2006}.
El acuerdo entre modelos parecidos no es evidencia independiente: si comparten variables y ventana de
entrenamiento, sus errores están correlacionados.

\subsection{Una referencia estructural: Martin y Wagner}\label{sec:martin}

\citet{martin2019} proponen una expresión del rendimiento esperado de una acción a partir de índices de
varianza implícita, derivada de \citet{martin2017}:
\begin{equation}
\frac{\E_t[R_{i,t+1}]-R_{f,t+1}}{R_{f,t+1}}=\mathrm{SVIX}_t^2+\tfrac12\Bigl(\mathrm{SVIX}_{i,t}^2-\overline{\mathrm{SVIX}}_t^2\Bigr),
\label{eq:martin}
\end{equation}
donde $\mathrm{SVIX}_t$ es el del mercado, $\mathrm{SVIX}_{i,t}$ el de la acción y
$\overline{\mathrm{SVIX}}_t^2$ el promedio ponderado por valor de la sección transversal. Queda como
referencia estructural si se dispone de esos insumos de mercado y de sección transversal; no sustituye la
calibración diaria descrita arriba.

\begin{noconcluir}
Ni la \eqref{eq:martin} ni ningún modelo de esta sección convierte a $\Q$ en $\P$. Todo pronóstico $\P$ se
valida con rendimientos realizados y con las reglas de puntuación de la \secref{sec:calibracion}.
\end{noconcluir}

\begin{criterio}
Mejora incremental fuera de muestra respecto de la referencia, dentro de la misma familia de modelos, o
una delimitación honesta de dónde no hay utilidad.
\end{criterio}
